\section{Simulation details}\label{app:sim}

\paragraph{Construction.} The person and firm effects estimated on the panel are the true effects. Every
scenario starts from baseline variances $\bar\sigma^2_{gt}$ for the person-years: a fit of the leave-out
estimates $\hat\sigma_i^2$ on the scale characteristics $z_i$ of Section~\ref{sec:transfer} alone (for stayers, on a
quadratic in the rank of the leverage), assigned to every person-year of the row. A scenario multiplies them
by a profile factor and, where stated, by a worker-specific factor. The factors are drawn once per scenario
and then held fixed, like the effects; in every draw only the errors are new, independent across person-years
with these variances, and the variance of $\hat\theta$ is computed for the fixed variances. Movers' rows are then formed as in Section~\ref{sec:setting}. Let $\tilde\alpha$
denote the true person effect standardised across person-years and $\tilde\omega$ the standardised sum of the
true person and firm effects.

\begin{table}[H]
\centering\small
\caption{Scenarios}\label{tab:scen}
\begin{tabular}{@{}rp{0.2\textwidth}p{0.7\textwidth}@{}}
\toprule
 & Errors & Variance of a person-year, relative to the baseline \\
\midrule
1 & normal & 1 \\
2 & normal & $1+0.5\,\tilde\alpha^2$ (rescaled to mean one), times a gamma factor per worker with mean one and coefficient of variation 1.5 \\
3 & symmetric, kurtosis 6 & as scenario 2 \\
4 & gamma, skewness 1 & $1+\tilde\alpha^2$ (rescaled); stayers' errors skewed against the wage level (left-skewed above the mean wage, right-skewed below) \\
5 & gamma, skewness 1 & as scenario 4, with stayers' errors skewed with the wage level \\
6 & normal & step function of $\tilde\alpha$: the stayers' sample variance in the data by decile of the person effect, relative to its mean; times the gamma factor of scenario 2 \\
7 & normal & $\exp(-0.6\,\tilde\omega)$ (rescaled) \\
8 & normal & stayers $1+0.5\,\tilde\alpha^2$, movers $1+1.5\,\tilde\alpha^2$ (each rescaled to mean one within its group) \\
9 & normal & stayers $1+\tilde\alpha^2$ (rescaled), movers 1 \\
10 & normal & stayers as in scenario 6; movers a step function of $\tilde\alpha$ on the same deciles, with the
values of the first row of Table~\ref{tab:deciles} relative to their mean (3.57, 1.13, 0.71, 0.58, 0.54, 0.53,
0.58, 0.73, 1.04, 1.85); times the gamma factor of scenario 2 \\
\bottomrule
\multicolumn{3}{@{}p{0.94\textwidth}@{}}{\rule{0pt}{2.6ex}Errors in scenario 3: a standardised gamma variable
with shape 2, with a random sign. In scenarios 4 and 5: a standardised gamma variable with shape 4 (skewness
1, kurtosis 4.5); for stayers its sign is fixed, opposite to the sign of the stayer's $m_i$ for the variance of
the sum (``against the wage level'') or equal to it; for movers the sign is random, so that movers' errors are
symmetric and the stayers' skewness is isolated.}
\end{tabular}
\end{table}

\paragraph{True variance.} For each scenario the benchmark for the variance of $\hat\theta$ is $4\sum_im_i^2\sigma_i^2$ plus
twice the double sum of \eqref{eq:V}, with the weights as used by the estimator, that is, \eqref{eq:ustat} with
the projected $B_{ii}$ and $P_{ii}$ on movers' rows; the double sum is computed as in \eqref{eq:sim} with 256 draws
at the true variances. On movers' rows these weights have a small diagonal that is not zero. The terms in third
and fourth cumulants that it adds to the variance in the scenarios without normal errors are not included, so
the benchmark is approximate there. These terms are of second order in the projection error of the diagonal
(a few percent with 1{,}600 projections); they were not computed. Across the scenarios the variance of the simulated estimates
is between \R{sim.mcvar.min} and \R{sim.mcvar.max} times this number, in line with the sampling error of a
variance from \R{sim.reps} draws (two standard errors are about 0.2), which makes this a coarse check only.

\paragraph{The model's error at fixed characteristics.} For Table~\ref{tab:bdelta} the two steps of
Section~\ref{sec:transfer} are run once per scenario with the true variances $\sigma_i^2$ as
outcomes and with the true effects, standardised and winsorised in the same way, in place of the estimated
ones. With $\Delta_i=v_i^*-\sigma_i^2$, the first part of $b(\Delta)$ in \eqref{eq:b} is computed directly. For the
double sum write $\Delta=a-b$ with $a$ and $b$ the positive and negative parts; then
$\sum\sum C_{i\ell}^2\Delta_i\Delta_\ell=2Q(a)+2Q(b)-Q(a+b)$ with $Q(v)=\sum\sum C_{i\ell}^2v_iv_\ell$, and each $Q$
is computed as in \eqref{eq:sim} with \R{sim.bd.draws} draws, the same for the three.

\paragraph{Realised scenarios.} Table~\ref{tab:regime} reports, for each scenario, the average variance factor
by decile of the true person effect, separately for stayers and movers, the movers' share of the first term of
\eqref{eq:V} for the variance of the sum, and the two parts of $b(\Delta)$ behind Table~\ref{tab:bdelta} for that
component.

\begin{table}[H]
\centering\footnotesize
\caption{Realised variance factors by decile of the person effect (lowest, fifth, highest); for
$\sigma^2_{\alpha+\psi}$, movers' share of the first term and the two parts of $b(\Delta)/\V[\hat\theta]$}\label{tab:regime}
\begin{tabular}{@{}rccccccccc@{}}
\toprule
 & \multicolumn{3}{c}{Stayers} & \multicolumn{3}{c}{Movers} & Movers' share & \multicolumn{2}{c}{$b(\Delta)/\V[\hat\theta]$} \\
\cmidrule(lr){2-4}\cmidrule(lr){5-7}\cmidrule(l){9-10}
Scenario & lowest & fifth & highest & lowest & fifth & highest & of term 1 & first part & double sum \\
\midrule
\tabinput{gen/tab_regime.tex}
\bottomrule
\end{tabular}
\end{table}

\paragraph{Validation runs.} The three runs of Section~\ref{sec:validation} use the scenario definitions of
Table~\ref{tab:scen}, with the baseline variances, the effects and the decile profile of scenario 6 computed on
each run's own panel, and the movers' multipliers of scenario 10 as in the table. Their error draws come from a
random stream that the earlier runs did not use. The code of the estimator, the list of scenarios, the number
of draws and the thresholds were fixed, with file hashes, before the first task ran; one task was run and
checked first, and the results of each run were then aggregated and compared with the thresholds by a script,
which lists every cell that fails. Table~\ref{tab:valcells} gives all cells.

\begin{table}[H]
\centering\footnotesize
\caption{Validation runs, every cell: ratio as in Table~\ref{tab:sim} and coverage of the 95\% interval}\label{tab:valcells}
\begin{tabular}{@{}rlcccccccc@{}}
\toprule
 & & \multicolumn{4}{c}{Ratio} & \multicolumn{4}{c}{Coverage} \\
\cmidrule(lr){3-6}\cmidrule(l){7-10}
 & Scenario & $\sigma^2_\alpha$ & $\sigma^2_\psi$ & $\sigma_{\alpha,\psi}$ & $\sigma^2_{\alpha+\psi}$ & $\sigma^2_\alpha$ & $\sigma^2_\psi$ & $\sigma_{\alpha,\psi}$ & $\sigma^2_{\alpha+\psi}$ \\
\midrule
\tabinput{gen/tab_val_cells.tex}
\bottomrule
\end{tabular}
\end{table}

\paragraph{Runs.} Eight tasks of 25 draws per scenario, on Boston University's Shared Computing Cluster; errors
are drawn from random streams keyed by scenario, task and draw, so results do not depend on the order of
execution. No draw was discarded. Two things were decided after first runs on these scenarios, and the reader
should know them. Scenario 10 was added when the comparison in Table~\ref{tab:deciles} had been seen. And the
form of the model for movers was chosen among four candidates that were run on the same draws. A rule fixed
before the comparison (the smallest worst shortfall among candidates with no ratio above 1.10 and with floors
binding for fewer than 0.5\% of movers' rows on the data, a condition that excluded one candidate) selected a
different candidate, which took the shape of the profile from the stayers and had ratios between 0.93 and
1.06. The present model exceeded the limit of that rule narrowly (1.103 and 1.104 in two cells). After seeing
the results we chose it all the same, because its errors are on the conservative side. The ratios in
Table~\ref{tab:sim} are therefore not an out-of-sample test of that choice; the validation of
Section~\ref{sec:validation} was run afterwards, with thresholds fixed in advance.
The right panels of Tables~\ref{tab:sim} and~\ref{tab:simcov} come from separate runs on the same draws of the
data (the simulated estimates agree to twelve digits). This
is a development study of a prototype, not a registered confirmation of an implementation.

\paragraph{The $q=1$ simulation (Section~\ref{sec:q1sim}).} \emph{Panels.} 100 firms in two groups of 50, with
log-normal sizes (the log of size is standard normal); 99 remain after pruning in the designs without a
bottleneck. 12,000 workers, each assigned to a group with
probability one half and observed in two to six consecutive years of a six-year window (bridging workers in
three to six). A worker's first firm is drawn with probability proportional to size among the firms of the
worker's group; in each later year the worker is assigned a newly drawn firm of the group with probability
0.10 (the draw can return the same firm, so actual changes are somewhat rarer). A bridging worker moves once, in a year drawn at random, to a firm of the other group and belongs to
that group from then on. Without a bottleneck all firms form one group. The panel is pruned to the
leave-one-out connected set, and rows are formed as in Section~\ref{sec:setting}.
\emph{Effects.} Firm effects are plus or minus half the gap, by group, plus a normal deviation with standard
deviation 0.12; person effects are a normal deviation with standard deviation 0.30 plus one half of the
centred effect of the firm of the worker's last observed year. \emph{Variances of a person-year.} Flat: 0.04. Profile:
$0.04\,(1+0.5\,\tilde\alpha^2)$, rescaled to mean 0.04, times a gamma factor per worker with mean one and
coefficient of variation one. Movers: $1+0.5\,\tilde\alpha^2$ for stayers and $1+1.5\,\tilde\alpha^2$ for movers,
each rescaled to mean one within its group, times 0.04 for stayers and 0.12 for movers. Bridge: 0.04, and 0.12
for the bridging workers. \emph{Errors.} Normal; in two designs a centred exponential variable (skewness 2,
kurtosis 9) scaled to the same variances. Effects and variances are fixed within a design; each draw has new
errors.

\emph{Estimates.} As in the other simulations: 1,600 projections for movers' leverages and weights, exact values
for stayers, the model for movers fitted in every draw with estimated effects in the profile, 16 draws for the
double sums, and the outcome centred at its person-year mean, which the propositions treat as a constant
(Section~\ref{sec:controls}). The power iteration stops when $\lambda_1$ changes by less than $10^{-9}$ in
relative terms, or after 200 steps. The true values of $\V[\hat\theta]$ and $\Sigma$ are those of the estimator
as implemented, given its projections: the weights include the centring and the projected leverages, which are
drawn once per design and held fixed across draws of the outcome. They are computed with the formulas for
normal errors, $4m'\Omega m+2\tr(C\Omega C\Omega)$ with $\Omega=\diag(\sigma_i^2)$ and its analogues for $\Sigma$,
the traces with 512 random draws; repeating that calculation with other draws in three designs changed
$\Sigma_{\rm rr}$ by at most \R{q1.bench.srr} and $\V[\hat\theta]$ by at most \R{q1.bench.v} in relative terms (standard
deviation over four repetitions; the largest value is for ten bridging workers). With normal errors these
benchmarks are exact up to that simulation error. With skewed errors they omit terms in the third and fourth
moments that are proportional to the diagonal of the weights, which is not exactly zero with projected
leverages and centring; there they are approximations. Coverage is measured against the true component and
does not depend on these benchmarks.

\emph{Runs.} The run reported in Section~\ref{sec:q1sim} has \R{q1.reps} draws for each of 27 designs, in tasks of
100 with random streams keyed by run, design, task and draw. The designs, the tolerances and the hashes of the
source files were written to a manifest before the run; the realised bottlenecks (Table~\ref{tab:q1cells}) were
checked before any outcome was drawn; and the results were aggregated only after every task had been checked
against the manifest. \R{q1.fail.reps} of \R{q1.draws.total} draws failed, and none was discarded. A
development run on other panels (\R{q1.dev.draws} draws per design) preceded it. There a first version of the
estimates, which used $\sum_iw_{i1}^2\hat\sigma_i^2$ for $\hat\Sigma_{bb}$ and the fitted value only when that number
was unusable, covered between 0.910 and 0.995, with two of 75 values below 0.92, both with ten bridging workers
and skewed errors; using the fitted value alone covered as little as 0.785 where the bridging workers are
noisier than the model says. (These are shares of 200 draws, with a Monte Carlo standard error of about 0.015,
so the two low values alone would not establish undercoverage.) The rule \eqref{eq:q1est} was adopted then, and
nothing was changed afterwards. In the run reported here the correlation of $\hat b_1$ and the remainder is at
most \R{q1.rho.true.max} in absolute value across designs, and $\hat\rho_{\rm lo}^2$ exceeded $2/3$ in at most
\R{q1.rholo.above} of the draws of a case (among draws in which both diagonal entries are positive), so the
condition of Corollary~\ref{cor:q1} is not put to a test.

\begin{table}[H]
\centering\footnotesize\setlength{\tabcolsep}{4.2pt}
\caption{The $q=1$ simulation, every design: curvature, coverage of the normal and of the $q=1$ interval for
$\hat\sigma^2_\alpha$, $\hat\sigma^2_\psi$ and $\hat\sigma_{\alpha,\psi}$, coverage of the normal interval for the
variance of the sum, share of draws in which the fitted number is the larger one in \eqref{eq:q1est}, and mean
length of the $q=1$ interval over mean length of the normal interval (on the draws in which both exist; the
ratio is formed for each component and the three ratios are averaged). The normal interval uses the 0.975
quantile of the normal distribution; a draw without an interval counts as not covered}\label{tab:q1cells}
\tabinput{gen/tab_q1cells.tex}
\end{table}

